% !TeX root = ../main.tex
BARCS fits time courses, covariate-adjusted effects, and factorial
interactions through a common beta-binomial regression interface.  The
library-total-conditional model specifies the denominator, and control
diagnostics assess calibration.  The sampling model determines how sequencing
depth and between-library variation enter the fit; the design matrix
determines which biological question the coefficient addresses.

In Cas13, adding the intermediate time point improved essential-gene recovery
modestly.  All four alternative methods ranked essential genes more strongly,
however, and aggregation-matched control scaling produced similar calibration
diagnostics across methods.  Because the input was normalized and
ComBat-corrected, this comparison provides limited evidence about the
underlying count likelihoods.  In IL2RA, donor-adjusted BARCS recovered more
validated regulators with fewer calls than the matched MAGeCK-MLE fit.
Waterbear and MAUDE recovered more members of the validation panel, and the
panel's limited coverage prevents interpreting recovery per call as positive
predictive value.

In the simulations, moderation improved genome-scale F1, although MAGeCK-MLE
retained slightly stronger ranking and a lower realized FDP.  Control-based
denominators improved interaction recovery in every simulated seed, with a
much larger gain in the one seed where the full-library fit collapsed.  These
denominator choices change the reference for abundance: full-library effects
describe changes in library share, whereas control-based effects describe
changes relative to the selected control class.  Non-targeting and safe-harbor
controls also differ in whether they absorb the shared cutting response.  The
appropriate control class depends on the effect being estimated and the
behavior of the controls.

The IL2RA cross-fitting analysis supports the guide-level operating point
examined here.  The null grid shows why that result cannot establish
gene-level calibration: correlated guides can inflate Stouffer gene
probabilities even when guide tests appear calibrated.  Scaling after matched
aggregation reduced this inflation but left appreciable error in some cells.
The beta-binomial fit also uses plug-in dispersion estimates; a residual-$t$
reference does not fully account for their uncertainty.  Nominal gene-level
FDR thresholds in these analyses therefore lack a general calibration
guarantee.  Further development should account for guide dependence in gene
aggregation and assess calibration using blocked permutations or independent
controls.

In the external audit, restoring full-library totals corrected the reported
Avana CB\textsuperscript{2} discovery count, but the residual null inflation
and poor DeWeirdt interaction recovery leave the broader calibration concern
intact. Preserving totals through filtering prevents an unintended change in
the model; it does not resolve the effects of dependence, processed counts, or
comparison-specific confounding.

BARCS is most directly interpretable with likelihood-compatible counts and
independent biological libraries.  Repeated harvests from one culture and FACS
bins partitioned from one donor require repeated-measures or joint models to
represent their dependence.  Donor indicators adjust mean differences but do
not supply that covariance structure.  For these designs, the fitted
coefficients and control diagnostics remain exploratory.  Confirmatory
applications require biological replication and calibration checks at the
guide or gene level used to report discoveries.
